\documentclass[12pt]{article}
\usepackage[utf8]{inputenc}
\usepackage[T1]{fontenc}
\usepackage{hyperref}
\title{Urak Lawoi' Language Resources for K-12 Curriculum Development}
\author{Compiled by Claude Sonnet 5 (claude-sonnet-5) for Andr\'as Kornai}
\date{}
\begin{document}
\maketitle

\section{Introduction}
Urak Lawoi' (ISO 639-3 urk) is an Austronesian (Malayic) language spoken by the Urak Lawoi' on the Andaman Sea islands of Thailand, including Lanta and Bulon. It is endangered and mostly oral. The usual script is Thai. Some Thai-script and romanised literacy materials are reported.

\section{Primary Educational Resources}
No school textbook was found in the search. The most relevant item is research on revitalisation and intercultural education:
\begin{itemize}
\item Mahidol University repository, ``Ethnic Identity, Language Revitalization, and Intercultural Education among the Urak Lawoi in Thailand''. It describes a Thai-based orthography and culturally grounded curricula on Lanta Island: \url{https://repository.li.mahidol.ac.th/handle/123456789/115863}
\end{itemize}

\section{Linguistic and Reference Materials}
\begin{itemize}
\item Hogan, D.\ W., with S.\ W.\ Pattemore (1988), \emph{Urak Lawoi': basic structures and a dictionary}, Pacific Linguistics C-109: \url{https://openresearch-repository.anu.edu.au/bitstream/1885/146628/1/PL-C109.pdf}
\item Wikipedia, Urak Lawoi' language: \url{https://en.wikipedia.org/wiki/Urak_Lawoi%CA%BC_language}
\item Wacana article (Universitas Indonesia): \url{https://scholarhub.ui.ac.id/wacana/vol10/iss1/7}
\item Hogan's later Lincom volume (1999) is mentioned in a book listing: \url{https://www.goodreads.com/book/show/5854440}
\end{itemize}

\section{Religious and Community Resources}
None found.

\section{Digital and Audio Resources}
\begin{itemize}
\item ``Urak Lawoi Language on Bulon Island: Documenting and Characterizing a Near-Vanishing Language'', a documentation study of 759 words in IPA: \url{https://so05.tci-thaijo.org/index.php/pimjournal/article/view/12028}
\item Endangered Languages Project entry: \url{https://www.endangeredlanguages.com/lang/1988?hl=pt-br}
\end{itemize}

\section{Curriculum Development Context}
Adults speak the language but children largely do not, so revitalisation is the priority. There is a published dictionary and grammar and an emerging community orthography and curriculum effort on Lanta. No formal school textbooks or Ministry materials were found.
\end{document}